\section{傅里叶反演公式}\label{sec:Approach}

\subsection{\texorpdfstring{$\delta$分布的逼近}{delta分布的逼近}}\label{sub:approach}
在\refsub{shannon}中，我们对带限函数提出了公式：如果$\mathrm{supp}(\mathcal{F} f)\subset[-\frac{T}{2},\frac{T}{2}]$，则有
\begin{equation*}
    f(t)=T\mathrm{sinc}(Tt)*f(t)
\end{equation*}

如果希望对更一般的函数$f(x)$，找到一个函数$g(x)$，使得$f(x)=g*f(x)$，自然会想到狄拉克$\delta$分布，然而它是一个奇异分布，我们只能从形式上定义一个$\delta$函数。事实上，$$!\exists g,\quad f*g=f,\quad\forall f$$
但很容易构造一个函数列$g_n$，使得$f*g_n\to f$\cite{folland}。回忆我们在\refsub{intuition_of_convolution}中提到的“平均化”卷积，设$g_a(x)=\frac{1}{2a}\chi_{[-a,a]}$，则$g_a*f$是$f$在区间$[x-a,x+a]$上的平均值，即$\frac{1}{2a}f*g_a(t)=\int_{x-a}^{x+a}f(x)\,dx$，若$f\in C(\mathbb{R})$，则$g_a*f\to f,a\to 0$。

一般地，我们考虑函数的缩放，设$g\in L^1(\mathbb{R})$，我们记
\[g_\lambda(x)=\frac{1}{\lambda}g\left(\frac{x}{\lambda}\right)\]
这类似于时域的伸缩在傅里叶变换下的表现，它是一种“保面积”的缩放，即：
\[\int_{-\infty}^{\infty}g_\lambda(x)\,dt=\int_{-\infty}^{\infty}g\left(\frac{y}{\lambda}\right)\,d\frac{y}{\lambda}=\int_{-\infty}^{\infty}g(y)\,dy\]
更一般地，我们给出接下来要频繁使用的引理：\begin{lemma}
    \[\int_{a}^{b}g_\lambda(x)\,dx=\int_{a/\lambda}^{b/\lambda}g(y)\,dy\]
\end{lemma}

\begin{theorem}
    设$g\in L^1(\mathbb{R}),\int_{-\infty}^{\infty}g(x)\,dx=1$，并记$\alpha=\int_{-\infty}^{0}g(x)\,dx,\beta=\int_{0}^{\infty}g(x)\,dx$。假设$f\in PC(\mathbb{R})$，则当$g(x)$时限或$f(x)$有界时，$f*g(x)$良定义，并且
    \[\lim_{\lambda\to 0}f*g_\lambda(x)=\alpha f(x+)+\beta f(x-)\]
    特别地，如果$f(x)$在点$x$连续，则有$\lim_{\lambda\to 0}f*g_\lambda(x)=f(x)$；如果$f(x)$在某个闭区间$[a,b]$上连续，则以上极限是一致收敛于$f(x)$的。
\end{theorem}
\begin{proof}
    以下证明中，我们需要使用$\epsilon-\delta$语言，暂时不用$\delta$表示狄拉克分布。
    \begin{align*}
        f*g_\lambda(x)-\alpha f(x+)-\beta f(x-)=\int_{-\infty}^{0}[f(x-y)-f(x+)]g_\lambda(y)\,dy+\int_{0}^{\infty}[f(x-y)-f(x-)]g_\lambda(y)\,dy
    \end{align*}
    我们希望证明在$\lambda$充分小时，以上两个积分也充分小。以$\int_{0}^{\infty}$为例，由于$f\in PC(\mathbb{R})$，存在$\delta>0$，使得当$|y|<\delta$时，$|f(x-y)-f(x-)|<\epsilon$。于是
    \begin{align*}
        \left|\int_{0}^{\delta}[f(x-y)-f(x-)]g_\lambda(y)\,dy\right|
        <\epsilon\left|\int_{0}^{\delta}g_\lambda(y)\,dy\right|=\epsilon\left|\int_{0}^{\delta/\lambda}g(y)\,dy\right|\leq\epsilon\int_{0}^{\infty}|g(y)|\,dy
    \end{align*}
    我们用剩下的条件估计$\int_{\delta}^{\infty}[f(x-y)-f(x-)]g_\lambda(y)\,dy$。如果$g(x)$时限，设$\supp(g)\subset[-T,T]$，则当$\lambda<\delta/T$时，$$y>\delta/\lambda>T\implies g(y)=0\implies\int_{\delta}^{\infty}[f(x-y)-f(x-)]g_\lambda(y)\,dy=\int_{\delta/\lambda}^{\infty}[f(x-\lambda y)-f(x-)]g(y)\,dy=0$$

    如果$f(x)$有界，设$|f(x)|<M$，则
    \begin{align*}
        \left|\int_{\delta}^{\infty}[f(x-y)-f(x-)]g_\lambda(y)\,dy\right| & \leq \int_{\delta}^{\infty}|f(x-y)-f(x-)||g_\lambda(y)|\,dy \\
                                                                          & \leq 2M\int_{\delta}^{\infty}|g_\lambda(y)|\,dy             \\
                                                                          & =2M\int_{\delta/\lambda}^{\infty}|g(y)|\,dy
    \end{align*}
    $\int_{-\infty}^{\infty}|g(y)|\,dy<\infty$的必要条件是$\int_{\delta/\lambda}^{\infty}|g(y)|\,dy\to 0,\delta/\lambda\to\infty$，因此上式随$\lambda\to 0$趋于0。

    最后，如果$f(x)$在闭区间$[a,b]$上连续，则它在该区间上一致连续，以上证明中的$\delta$与$x$无关，极限是一致收敛的。
\end{proof}
一种常用的函数列是高斯函数列，这里取一个积分值为1的特例：
\[G(x)=\frac{1}{\sqrt{\pi}}\mathrm{e}^{-x^2},\quad G_\lambda(x)=\frac{1}{\lambda\sqrt{\pi}}\mathrm{e}^{-x^2/\lambda^2}\]
尽管用很多其他函数都可以逼近$\delta$分布，但高斯函数列有一系列非常好的性质：\begin{itemize}
    \item 高斯函数是偶函数，$\alpha=\beta=\dfrac{1}{2}$；
    \item 高斯函数是无限阶可导的。我们知道，卷积能够继承函数的可导性：$f*g'=f'*g=(f*g)'$，使用高斯函数做卷积，所得的函数光滑
    \item 高斯函数是施瓦兹函数，即$G_\lambda\in \mathcal{S}$
    \item 高斯函数的傅里叶变换仍是高斯函数，这使得傅里叶变换更容易计算
    \item 高斯函数$G_\lambda\in L^1(\mathbb{R})$并且是有界的，与任意$f\in PC(\mathbb{R})$卷积时，卷积积分良定义且趋向于$f$
\end{itemize}

作为以上定理的推论，可以得到多项式空间在连续函数空间中的稠密性：
\begin{theorem}[*魏尔斯特拉斯第一逼近定理]\label{thm:weierstrass_1}
    设$f\in C([a,b]),-\infty<a<b<\infty$，则$f$是某个多项式序列在$[a,b]$上的一致极限，即：
    \[\forall \epsilon>0,\exists \text{多项式}P,\quad \sup_{a\leq x\leq b}|f(x)-P(x)|<\epsilon\]
\end{theorem}
%在处理分析问题时，类似分段光滑函数可由三角级数一致逼近，将连续函数用多项式一致逼近是一个很强大的工具。
\begin{proof}
    将$f$做连续延拓，使得
    $$f\in C(\mathbb{R}),\quad f(x)=0,\quad\forall x\in (-\infty,a-1]\cup[b+1,\infty)$$
    则$\lim_{\lambda\to 0}f*G_\lambda \rightrightarrows f,\forall x\in[a,b]$，即$f*G_\lambda$在区间$[a,b]$上一致收敛到$f$。具体来说，
    \[\forall \epsilon>0,\exists \lambda>0,\quad \sup_{a\leq x\leq b}\left|f(x)-\frac{1}{\lambda\sqrt{\pi}}\int_{a-1}^{b+1}\mathrm{e}^{-(x-y)^2/\lambda^2}f(y)\,dy\right|<\frac{\epsilon}{2}\]
    考虑将右边的积分用多项式近似。我们知道，$\mathrm{e}^{-t^2}$可以展开为泰勒级数，并且在$\mathbb{C}$上内闭一致收敛，即
    \[\sum_{n=0}^{N}\frac{t^n}{n!}\rightrightarrows \mathrm{e}^{-t^2},\quad N\to\infty\]
    具体来说，取$|f(x)|$的上界$M$（紧支连续函数是一致连续的，必有上界），存在$N$，使得
    \[\sup_{t\in[a-1,b+1]}\left|\mathrm{e}^{-t^2}-\sum_{n=0}^{N}\frac{(-1)^n t^{2n}}{n!}\right|<\frac{\epsilon\sqrt{\pi}}{2M(b-a+2)}\]
    于是
    \[P(x)=\frac{1}{\lambda\sqrt{\pi}}\sum_{0}^{N}\int_{a-1}^{b+1}\frac{(-1)^n (x-y)^{2n}}{\lambda^{2n}n!}f(y)\,dy,\quad \sup_{a\leq x\leq b}\left|f(x)-P(x)\right|<\epsilon\]
    我们还可以求出$P$的显式表达式：
    \[P(x)=\sum_{n=0}^{N}\sum_{k=0}^{2n}c_{kn}x^k,\quad c_{kn}=\frac{(-1)^{k-n}}{\lambda^{2n+1}n!\sqrt{\pi}}\binom{2n}{k}\int_{a-1}^{b+1}y^{2n-k}f(y)\,dy\]
    使用其他光滑函数列逼近$\delta$分布也可以得到不同的多项式逼近，不过，绝大多数情况下我们并不需要计算出这个表达式，而是直接用多项式的一致逼近来获取函数的性质。
\end{proof}

进一步，可以得到三角多项式（\cref{def:trinomial}）空间在连续的周期函数空间中的稠密性：
\begin{corollary}[魏尔斯特拉斯第二逼近定理]\label{cor:weierstrass_2}
    设$f\in C(\mathbb{R})$以$T$为周期，则$f$是某个同周期的三角多项式序列的一致极限，即：
    \[\forall \epsilon>0,\exists \text{三角多项式}P,\quad \sup_{x}|f(x)-P(x)|<\epsilon\]
\end{corollary}
\begin{proof}
    由于可以对自变量进行伸缩变换，我们不妨设$T=2\pi$，并重点考虑区间$[-\pi,\pi]$.

    首先，三角多项式的奇分量和偶分量是显式地分离的，这启发我们考虑将$f(t)$也分解为奇分量和偶分量：
    $$f_\e(t)=\frac{f(t)+f(-t)}{2},\quad f_\text{o}(t)=\frac{f(t)-f(-t)}{2}$$
    一种直接的想法是，通过变量代换$x=\cos t$和$x=\sin t$，分别将逼近$f_e(\arccos x)$和$f_\text{o}(\arcsin x)$的多项式转化为三角多项式。然而，$\sin(k x),k\in\mathbb{N}$有共同的零点和对称轴：
    $$\sin(-k\pi)=\sin 0=\sin(k\pi)=0,\quad \sin(k(\pi-x))=\sin(kx)=\sin(k(-\pi-x))$$
    对于奇分量来说，我们只能保证
    $$f_\text{o}(-\pi)=f_\text{o}(0)=f_\text{o}(\pi)=0$$
    而不能保证$f_\text{o}(t)$关于$t=\pm \dfrac{\pi}{2}$轴对称，因此不能单独使用正弦多项式逼近奇分量。

    考虑以下两个偶函数：
    $$g(t)=f(t)+f(-t),\quad h(t)=\sin(t)(f(t)-f(-t))$$
    做变量代换$x=\cos t$，利用\cref{thm:weierstrass_1}，有：
    $$\forall \epsilon>0,\exists\text{多项式}P,Q,\quad \sup_{-1\le x\le 1}|g(\arccos x)-P(x)|<\epsilon,\sup_{-1\le x\le 1}|h(\arccos x)-Q(x)|<\epsilon$$
    代回变量$t$，$P(\cos t)$和$Q(\cos t)$都是以$2\pi$为周期的偶函数，符合我们的要求。

    我们希望能用三角多项式逼近$f(t)-f(-t)$，也就是说$Q(\cos t)/\sin t$仍为三角多项式，这实际上要求$Q(\cos t)$在$t=0$和$\pm \pi$处，即$x=\pm 1$处有零点。注意到$x=\pm 1$时，有：
    $$h(\arccos x)=0,\quad |Q(x)-h(\arccos x)|=|Q(x)|<\epsilon$$
    构造修正多项式
    $$R(x)=Q(x)-Q(1)\frac{1+x}{2}-Q(-1)\frac{1-x}{2}$$
    这样，$R(\pm 1)=0$，且
    \begin{align*}
        |R(x)-Q(x)| & =\left|Q(1)\frac{1+x}{2}+Q(-1)\frac{1-x}{2}\right| \\
                    & \le|Q(1)|\frac{|1+x|}{2}+|Q(-1)|\frac{|1-x|}{2}    \\
                    & \le 2\epsilon,\quad \forall x\in[-1,1]
    \end{align*}
    因此
    $$\sup_{-1\le x\le 1}|R(x)-h(\arccos x)|\le\sup_{-1\le x\le 1}|R(x)-Q(x)|+\sup_{-1\le x\le 1}|Q(x)-h(\arccos x)|\le 3\epsilon$$
    由于$R(x)$有零点$\pm 1$，它一定有因式$(1-x)(1+x)=1-x^2$，记$R(x)=R_0(x)(1-x^2)$，则
    $$R(\cos t)=R_0(\cos t)(1-\cos^2 t)=R_0(\cos t)\sin^2 t$$

    最后，我们通过分段估计来说明三角多项式$R_0(\cos t)\sin t$能够一致逼近$F(t)=f(t)-f(-t)$.一方面，由于$F(t)$连续且$F(0)=F(\pi)=0$（根据周期性，$F(-\pi)$处的情况与$F(\pi)$处相同），有
    $$\forall\eta>0,\exists\delta_1>0,|t|<\delta_1\text{或}|t-\pi|<\delta_1\implies |F(t)|<\frac{\eta}{2}$$
    取$R_0(\cos t)$的上界$M$，根据$\sin t$的连续性，
    $$\exists\delta_2>0,|t|<\delta_2\text{或}|t-\pi|<\delta_2\implies|\sin t|<\frac{\eta}{2M}$$
    令$\delta=\min\{\delta_1,\delta_2\}$，则
    $$|t|<\delta\text{或}|t-\pi|<\delta\implies |F(t)-R_0(\cos t)\sin t|\le|F(t)|+M|\sin t|\le\eta$$
    另一方面，当$|t|>\delta$且$|t-\pi|>\delta$时，$|sin t|>\sin\delta>0$，因此
    $$|F(t)\sin t-R_0(\cos t)\sin^2 t|<\epsilon\implies|F(t)-R_0(\cos t)\sin t|<\frac{\epsilon}{\sin\delta}$$
    只要选取$\epsilon<\eta\sin\delta$，就能保证$\sup_t |F(t)-R_0(\cos t)\sin t|<\eta$.
\end{proof}

\subsection{黎曼-勒贝格引理}\label{sub:riemman_lebesgue_lemma}
本小节来补充\refsub{introduction_of_fourier}中提到的\textbf{黎曼-勒贝格引理}的证明。
\begin{theorem}[黎曼-勒贝格引理]\label{thm:riemann_lebesgue}
    设$f\in L^1(\mathbb{R})$，则其傅里叶变换$\mathcal{F} f(x)\to 0,x\to\infty$。
\end{theorem}
我们给出一个基于阶梯函数逼近的证明，尽管现在不需要掌握它，但可以从这个证明体会到分析学中逼近技术的强大威力。
\begin{proof}
    首先假设$f$是阶梯函数，即$f(t)=\sum_{k=1}^{n}c_k \chi_{[a_k,b_k]}(t)$，我们知道：\lr{
    \mathcal{F} \chi_{[a,b]}(\xi)=\int_{a}^{b}\mathrm{e}^{-2\pi \mathrm{i}\xi t}\,dt=\frac{\mathrm{e}^{-2\pi \mathrm{i}\xi a}-\mathrm{e}^{-2\pi \mathrm{i}\xi b}}{2\pi \mathrm{i}\xi}
    }{
    \mathcal{F} \chi_{[a,b]}(\omega)=\int_{a}^{b}\mathrm{e}^{-\mathrm{i}\omega t}\,dt=\frac{\mathrm{e}^{-\mathrm{i}\omega a}-\mathrm{e}^{-\mathrm{i}\omega b}}{\mathrm{i}\omega}
    }
    分子是两个模值为1的复数相减，分母趋于无穷，故
    $$\mathcal{F} \chi_{[a,b]}(x)\to 0,\quad x\to\infty$$
    由线性性，得到：
    $$\mathcal{F} f(x)=\sum_{k=1}^n c_k\mathcal{F}\chi_{[a_k,b_k]}\to 0,\quad x\to\infty$$

    设$f\in L^1(\mathbb{R})$，则根据实变函数的理论，存在阶梯函数列$f_n$，使得$\|f-f_n\|_1\to 0,n\to\infty$。我们知道，
    \[|\mathcal{F} g(x)|\leq \int_{-\infty}^{\infty}|g(t)|\,dt=\|g\|_1<\infty,\quad g\in L^1(\mathbb{R})\]
    因此，
    \[\sup_{x\in\mathbb{R}}|\mathcal{F} f(x)-\mathcal{F} f_n(x)|\leq \|f-f_n\|_1\to 0,\quad n\to\infty\]
    即$\mathcal{F} f_n\rightrightarrows\mathcal{F} f,\forall x\in\mathbb{R}$，于是$\mathcal{F} f(x)\to 0,x\to\infty$，得证。
\end{proof}

\subsection{傅里叶反演公式}\label{sub:fourier_inversion}
\begin{theorem}[傅里叶反演公式]
    \label{thm:fourier_inversion_1}
    设$f\in L^1(\mathbb{R})\cap PC(\mathbb{R})$，则对于任意点$x$，都有
    \lr{
    &\frac{f(x+)+f(x-)}{2}\\
    =&\lim_{\lambda\to 0}\int_{-\infty}^{\infty}\mathrm{e}^{2\pi \mathrm{i}\xi x}\mathcal{F} f(\xi)\mathrm{e}^{-\pi^2 \lambda^2 \xi^2}\,d\xi
    }{
    &\frac{f(x+)+f(x-)}{2}\\
    =&\lim_{\lambda\to 0}\frac{1}{2\pi}\int_{-\infty}^{\infty}\mathrm{e}^{\mathrm{i}\omega x}\mathcal{F} f(\omega)\mathrm{e}^{-\lambda^2\omega^2/4}\,d\omega
    }
    如果还有$\mathcal{F} f(x)\in L^1(\mathbb{R})$，则有
    \lr{
    \frac{f(x+)+f(x-)}{2}=\int_{-\infty}^{\infty}\mathrm{e}^{2\pi \mathrm{i}\xi x}\mathcal{F} f(\xi)\,d\xi
    }{
    \frac{f(x+)+f(x-)}{2}=\frac{1}{2\pi}\int_{-\infty}^{\infty}\mathrm{e}^{\mathrm{i}\omega x}\mathcal{F} f(\omega)\,d\omega
    }
\end{theorem}
我们指出，$\mathcal{S} \subset L^1(\mathbb{R})\cap PC(\mathbb{R})$。

要说明傅里叶逆变换的确能够还原出时域函数，一种直接的想法是积分换序，即\lr{
&\int_{-\infty}^{\infty}\mathrm{e}^{2\pi \mathrm{i}\xi x}\mathcal{F} f(\xi)\,d\xi\\
=&\int_{-\infty}^{\infty}\left(\int_{-\infty}^{\infty}f(x)\mathrm{e}^{-2\pi \mathrm{i}\xi x}\,dx\right)\mathrm{e}^{2\pi \mathrm{i}\xi x}\,d\xi
}{
&\frac{1}{2\pi}\int_{-\infty}^{\infty}\mathrm{e}^{\mathrm{i}\omega x}\mathcal{F} f(\omega)\,d\omega\\
=&\frac{1}{2\pi}\int_{-\infty}^{\infty}\left(\int_{-\infty}^{\infty}f(x)\mathrm{e}^{-\mathrm{i}\omega x}\,dx\right)\mathrm{e}^{\mathrm{i}\omega x}\,d\omega
}
然而，\lr{
\int_{-\infty}^{\infty}\mathrm{e}^{2\pi \mathrm{i}\xi (x-y)}\,d\xi
}{
\int_{-\infty}^{\infty}\mathrm{e}^{\mathrm{i}\omega (x-y)}\,d\omega
}
两个积分是发散的，无法进行积分换序。我们可以使用“收敛因子”来处理这个问题，从信号处理的角度，这就是\textbf{加窗}(windowing)的手法，见\refsub{window}。这里我们用高斯窗，仍记$G_\lambda(\xi)=\frac{1}{\lambda\sqrt{\pi}}\mathrm{e}^{-\xi^2/\lambda^2}$，有\lr{
    \mathcal{F} G_\lambda(\xi)=\mathrm{e}^{-\pi^2 \lambda^2 \xi^2}
}{
    \mathcal{F} G_\lambda(\omega)=\mathrm{e}^{-\lambda^2 \omega^2/4}
}
下面的证明中将用到高斯函数的傅里叶逆变换，但我们可以不依赖于傅里叶反演公式而预先求出它，具体的做法可以参考\cref{exmp:gauss_fourier}中求傅里叶正变换的过程。
\begin{proof}
    考虑
    \lr{
    &\int_{-\infty}^{\infty}\mathrm{e}^{2\pi \mathrm{i}\xi x}\mathcal{F} f(\xi)\mathrm{e}^{-\pi^2\lambda^2 \xi^2}\,d\xi\\
    =&\int_{-\infty}^{\infty}\left(\int_{-\infty}^{\infty}f(x)\mathrm{e}^{-2\pi \mathrm{i}\xi x}\,dx\right)\mathrm{e}^{2\pi \mathrm{i}\xi x}\mathrm{e}^{-\pi^2 \lambda^2 \xi^2}\,d\xi\\
    =&\int_{-\infty}^{\infty}f(x)\,dx\int_{-\infty}^{\infty}\mathrm{e}^{2\pi \mathrm{i}\xi (x-y)}\mathrm{e}^{-\pi^2 \lambda^2 \xi^2}\,d\xi
    }{
    &\frac{1}{2\pi}\int_{-\infty}^{\infty}\mathrm{e}^{\mathrm{i}\omega x}\mathcal{F} f(\omega)\mathrm{e}^{-\lambda^2\omega^2/4}\,d\omega\\
    =&\frac{1}{2\pi}\int_{-\infty}^{\infty}\left(\int_{-\infty}^{\infty}f(x)\mathrm{e}^{-\mathrm{i}\omega x}\,dx\right)\mathrm{e}^{\mathrm{i}\omega x}\mathrm{e}^{-\lambda^2\omega^2/4}\,d\omega\\
    =&\frac{1}{2\pi}\int_{-\infty}^{\infty}f(x)\,dx\int_{-\infty}^{\infty}\mathrm{e}^{\mathrm{i}\omega (x-y)}\mathrm{e}^{-\lambda^2\omega^2/4}\,d\omega
    }

    以上过程利用了高斯函数的速降性质，保证了积分的绝对收敛性，从而可以进行积分换序。内层的积分恰为傅里叶变换：
    \lr{
    &\int_{-\infty}^{\infty}\mathrm{e}^{2\pi \mathrm{i}\xi (x-y)}\mathrm{e}^{-\pi^2 \lambda^2 \xi^2}\,d\xi\\
    =&\mathcal{F}^{-1}\left[\mathrm{e}^{-\pi^2 \lambda^2 \xi^2}\right](x-y)\\
    =&\frac{1}{\sqrt{\pi}\lambda}\mathrm{e}^{-(x-y)^2/\lambda^2}\\
    =&G_{\lambda}(x-y)
    }{
    &\frac{1}{2\pi}\int_{-\infty}^{\infty}\mathrm{e}^{\mathrm{i}\omega (x-y)}\mathrm{e}^{-\lambda^2\omega^2/4}\,d\omega\\
    =&\mathcal{F}^{-1}\left[\mathrm{e}^{-\lambda^2\omega^2/4}\right](x-y)\\
    =&\frac{1}{\sqrt{\pi}\lambda}\mathrm{e}^{-(x-y)^2/\lambda^2}\\
    =&G_{\lambda}(x-y)
    }
    于是我们得到
    \lr{
    &\int_{-\infty}^{\infty}\mathrm{e}^{2\pi \mathrm{i}\xi x}\mathcal{F} f(\xi)\mathrm{e}^{-\pi^2 \lambda^2 \xi^2}\,d\xi\\
    =&\int_{-\infty}^{\infty}f(y)G_{\lambda}(x-y)\,dy\\
    =&f*G_{\lambda}(x)
    }{
    &\frac{1}{2\pi}\int_{-\infty}^{\infty}\mathrm{e}^{\mathrm{i}\omega x}\mathcal{F} f(\omega)\mathrm{e}^{-\lambda^2\omega^2/4}\,d\omega\\
    =&\int_{-\infty}^{\infty}f(y)G_{\lambda}(x-y)\,dy\\
    =&f*G_{\lambda}(x)
    }
    我们知道，$\lim_{\lambda\to 0}f*G_{\lambda}(x)=\frac{f(x+)+f(x-)}{2}$，因此
    \begin{lralign}
        &\lim_{\lambda\to 0}\int_{-\infty}^{\infty}\mathrm{e}^{2\pi \mathrm{i}\xi x}\mathcal{F} f(\xi)\mathrm{e}^{-\pi^2 \lambda^2 \xi^2}\,d\xi\\
        =&\frac{f(x+)+f(x-)}{2}
        \switch
        &\lim_{\lambda\to 0}\frac{1}{2\pi}\int_{-\infty}^{\infty}\mathrm{e}^{\mathrm{i}\omega x}\mathcal{F} f(\omega)\mathrm{e}^{-\lambda^2\omega^2/4}\,d\omega\\
        =&\frac{f(x+)+f(x-)}{2}
    \end{lralign}
    现在，只需要证明$\lim_{\lambda\to 0}$与积分号可以交换顺序。我们有：\lr{
    \left|\mathrm{e}^{2\pi \mathrm{i}\xi x}\mathcal{F} f(\xi)\mathrm{e}^{-\pi^2 \lambda^2 \xi^2}\right|\leq & |\mathcal{F} f(\xi)|\in L^1(\mathbb{R})
    }{
    \left|\frac{1}{2\pi}\mathrm{e}^{\mathrm{i}\omega x}\mathcal{F} f(\omega)\mathrm{e}^{-\lambda^2\omega^2/4}\right|\leq & \frac{1}{2\pi}|\mathcal{F} f(\omega)|\in L^1(\mathbb{R})
    }
    因此可以使用控制收敛定理，得到
    \lr{
    &\frac{f(x+)+f(x-)}{2}\\
    =&\lim_{\lambda\to 0}\int_{-\infty}^{\infty}\mathrm{e}^{2\pi \mathrm{i}\xi x}\mathcal{F} f(\xi)\mathrm{e}^{-\pi^2 \lambda^2 \xi^2}\,d\xi\\
    =&\int_{-\infty}^{\infty}\mathrm{e}^{2\pi \mathrm{i}\xi x}\mathcal{F} f(\xi)\,d\xi
    }{
    &\frac{f(x+)+f(x-)}{2}\\
    =&\lim_{\lambda\to 0}\frac{1}{2\pi}\int_{-\infty}^{\infty}\mathrm{e}^{\mathrm{i}\omega x}\mathcal{F} f(\omega)\mathrm{e}^{-\lambda^2\omega^2/4}\,d\omega\\
    =&\frac{1}{2\pi}\int_{-\infty}^{\infty}\mathrm{e}^{\mathrm{i}\omega x}\mathcal{F} f(\omega)\,d\omega
    }
\end{proof}
傅里叶反演公式还有一些其他的表述。例如，如果我们将窗函数换为矩形窗，则可得到：
\begin{theorem}[傅里叶反演公式]
    \label{thm:fourier_inversion}
    设$f\in L^1(\mathbb{R})\cap PS(\mathbb{R})$，则对于任意点$x$，都有
    \lr{
    &\frac{f(x+)+f(x-)}{2}\\
    =&\lim_{R\to \infty}\int_{-R}^{R}\mathrm{e}^{2\pi \mathrm{i}\xi x}\mathcal{F} f(\xi)\,d\xi
    }{
    &\frac{f(x+)+f(x-)}{2}\\
    =&\lim_{R\to \infty}\frac{1}{2\pi}\int_{-R}^{R}\mathrm{e}^{\mathrm{i}\omega x}\mathcal{F} f(\omega)\,d\omega
    }
\end{theorem}
等式右侧的积分与反常积分是有区别的，无穷限的反常积分要求积分上下限以二重极限的方式趋于无穷，而这里给出的是\textbf{主值意义下的积分}，即
\[\int_{-\infty}^{\infty}f(x)\,dx=\lim_{a\to-\infty,b\to\infty}\int_{a}^{b}f(x)\,dx\neq \lim_{R\to\infty}\int_{-R}^{R}f(x)\,dx=\mathrm{p.v.}\int_{-\infty}^{\infty}f(x)\,dx\]
\begin{example}
    设$f(x)=x$，则$\int_{-\infty}^{\infty}f(x)\,dx$发散，但$\mathrm{p.v.}\int_{-\infty}^{\infty}f(x)\,dx=0$
\end{example}
\begin{proof}
    我们有\lr{
    \int_{-R}^{R}\mathrm{e}^{2\pi \mathrm{i}\xi x}\,d\xi=&\frac{\mathrm{e}^{2\pi \mathrm{i} R x}-\mathrm{e}^{-2\pi \mathrm{i} R x}}{2\pi \mathrm{i} x}\\
    =&\frac{\sin(2\pi R x)}{\pi x}
    }{
    \frac{1}{2\pi}\int_{-R}^{R}\mathrm{e}^{\mathrm{i}\omega x}\,d\omega=&\frac{\mathrm{e}^{\mathrm{i} R x}-\mathrm{e}^{-\mathrm{i} R x}}{2\mathrm{i} x}\\
    =&\frac{\sin(R x)}{\pi x}
    }
    以及\cref{thm:dirichlet_int}：
    \[\int_{-\infty}^{0}\frac{\sin(ax)}{\pi x}\,dx=\int_{0}^{\infty}\frac{\sin(ax)}{\pi x}\,dx=\frac{1}{2}\]
    因此
    \begin{lralign}
        &\int_{-R}^{R}\mathrm{e}^{2\pi \mathrm{i}\xi x}\mathcal{F} f(\xi)\,d\xi\\
        =&\int_{-\infty}^{\infty}f(y)\,dy\int_{-R}^{R}\mathrm{e}^{2\pi \mathrm{i}\xi (x-y)}\,d\xi\\
        =&\int_{-\infty}^{\infty}f(y)\frac{\sin(2\pi R (x-y))}{\pi (x-y)}\,dy\\
        =&f*\left(\frac{\sin(2\pi R x)}{\pi x}\right)
        \switch
        &\frac{1}{2\pi}\int_{-R}^{R}\mathrm{e}^{\mathrm{i}\omega x}\mathcal{F} f(\omega)\,d\omega\\
        =&\int_{-\infty}^{\infty}f(y)\,dy\frac{1}{2\pi}\int_{-R}^{R}\mathrm{e}^{\mathrm{i}\omega (x-y)}\,d\omega\\
        =&\int_{-\infty}^{\infty}f(y)\frac{\sin(R (x-y))}{\pi (x-y)}\,dy\\
        =&f*\left(\frac{\sin(R x)}{\pi x}\right)
    \end{lralign}
    但是，由于$\mathrm{sinc}\notin L^1(\mathbb{R})$，无法保证$\lim_{R\to \infty} f(x)*\mathrm{sinc}(R x)=f(x)$，因此需要条件$f\in L^1(\mathbb{R})\cup PS(\mathbb{R})$，做更细致的估计：
    \lr{
    &\int_{-R}^{R}\mathrm{e}^{2\pi \mathrm{i}\xi x}\mathcal{F} f(\xi)\,d\xi-\frac{f(x+)+f(x-)}{2}\\
    =&\int_{-\infty}^{\infty}f(y)\frac{\sin(2\pi R (x-y))}{\pi (x-y)}\,dy-\frac{f(x+)+f(x-)}{2}\\
    =&\int_{-\infty}^{\infty}f(x-y)\frac{\sin(2\pi R y)}{\pi y}\,dy-\frac{f(x+)+f(x-)}{2}\\
    =&\int_{-\infty}^{0}[f(x-y)-f(x+)]\frac{\sin(2\pi R y)}{\pi y}\,dy\\
    &+\int_{0}^{\infty}[f(x-y)-f(x-)]\frac{\sin(2\pi R y)}{\pi y}\,dy
    }{
    &\frac{1}{2\pi}\int_{-R}^{R}\mathrm{e}^{\mathrm{i}\omega x}\mathcal{F} f(\omega)\,d\omega-\frac{f(x+)+f(x-)}{2}\\
    =&\int_{-\infty}^{\infty}f(y)\frac{\sin(R (x-y))}{\pi (x-y)}\,dy-\frac{f(x+)+f(x-)}{2}\\
    =&\int_{-\infty}^{\infty}f(x-y)\frac{\sin(R y)}{\pi y}\,dy-\frac{f(x+)+f(x-)}{2}\\
    =&\int_{-\infty}^{0}[f(x-y)-f(x+)]\frac{\sin(R y)}{\pi y}\,dy\\
    &+\int_{0}^{\infty}[f(x-y)-f(x-)]\frac{\sin(R y)}{\pi y}\,dy
    }
    我们需要证明当$R\to\infty$时，以上两个积分都趋于0。以$\int_{0}^{\infty}$为例，仍然考虑分段估计：$\int_{0}^{\infty}=\int_{0}^{K}+\int_{K}^{\infty}$，取$K\geq 1$，有
    \lr{
        &\left|\int_{K}^{\infty}\frac{\sin(2\pi Ry)}{\pi y}f(x-y)\,dy\right|\\
        \leq & \int_{K}^{\infty}\left|\frac{\sin(2\pi Ry)}{\pi y}\right| |f(x-y)|\,dy\\
        \leq & \int_{K}^{\infty}|f(x-y)|\,dy\\
        &\int_{K}^{\infty}\frac{\sin(2\pi Ry)}{\pi y}f(x-)\,dy\\
        = & f(x-)\int_{RK}^{\infty}\frac{\sin(2\pi Ry)}{\pi y}\,dy
    }{
        &\left|\int_{K}^{\infty}\frac{\sin(R y)}{\pi y}f(x-y)\,dy\right|\\
        \leq & \int_{K}^{\infty}\left|\frac{\sin(R y)}{\pi y}\right| |f(x-y)|\,dy\\
        \leq & \int_{K}^{\infty}|f(x-y)|\,dy\\
        &\int_{K}^{\infty}\frac{\sin(R y)}{\pi y}f(x-)\,dy\\
        = & f(x-)\int_{RK}^{\infty}\frac{\sin(R y)}{\pi y}\,dy
    }
    取$K$充分大，则以上两积分都是收敛的无穷限积分的尾部，从而是趋于0的。

    对于$\int_{0}^{K}$部分，定义\begin{align*}
        g(y)=\begin{cases}
                 \dfrac{f(x-y)-f(x-)}{\pi y}, & y\in [0,K]         \\
                 0,                           & \mathrm{otherwise}
             \end{cases}
    \end{align*}
    则\begin{lralign}
        &\int_{0}^{K}[f(x-y)-f(x-)]\frac{\sin(2\pi R y)}{\pi y}\,dy\\
        &=\int_{0}^{K}g(y)\frac{\mathrm{e}^{2\pi \mathrm{i} Ry}-\mathrm{e}^{-2\pi \im R y}}{2\mathrm{i}}\,dy\\
        &=\frac{1}{2\mathrm{i}}[\mathcal{F} g(-R)-\mathcal{F} g(R)]
        \switch
        &\int_{0}^{K}[f(x-y)-f(x-)]\frac{\sin(R y)}{\pi y}\,dy\\
        &=\int_{0}^{K}g(y)\frac{\mathrm{e}^{\mathrm{i} R y}-\mathrm{e}^{-\mathrm{i} R y}}{2\mathrm{i}}\,dy\\
        &=\frac{1}{2\mathrm{i}}[\mathcal{F} g(-R)-\mathcal{F} g(R)]
    \end{lralign}
    由于$f\in PS(\mathbb{R}),g\in PC(\mathbb{R}^*),g(0+)=f'(x-)$，有$g\in L^1(\mathbb{R})$，根据黎曼-勒贝格引理（\cref{thm:riemann_lebesgue}），$\mathcal{F} g(\pm R)\to 0,R\to\infty$，故$$\int_{0}^{K}[f(x-y)-f(x-)]\dfrac{\sin(R y)}{\pi y}\,dy\to 0,\quad R\to\infty$$
    得证。
\end{proof}

此外，还有一种类似的的傅里叶反演公式，我们仅仅给出其表述\cite{folland}：
\begin{claim}[傅里叶反演公式]
    设$f\in L^1(\mathbb{R})$，则
    \lr{
    &\frac{f(x+)+f(x-)}{2}\\
    =&\lim_{\lambda\to 0}\int_{-\infty}^{\infty}\mathrm{e}^{2\pi \mathrm{i}\xi x}\mathcal{F} f(\xi)\mathrm{e}^{-\pi^2 \lambda^2 \xi^2}\,d\xi,\quad\mathrm{a.e.}
    }{
    &\frac{f(x+)+f(x-)}{2}\\
    =&\lim_{\lambda\to 0}\frac{1}{2\pi}\int_{-\infty}^{\infty}\mathrm{e}^{\mathrm{i}\omega x}\mathcal{F} f(\omega)\mathrm{e}^{-\lambda^2\omega^2/4}\,d\omega,\quad\mathrm{a.e.}
    }
\end{claim}
